%% Manuscript: Site Effects in Multi-Centre EEG Data for Neuroscreening
%% Target journal: Biomedical Signal Processing and Control (Elsevier)
%% LaTeX class: elsarticle, numerical citations
%%
%% Compile with:
%%   pdflatex manuscript_eeg_harmonization.tex
%%   bibtex   manuscript_eeg_harmonization
%%   pdflatex manuscript_eeg_harmonization.tex
%%   pdflatex manuscript_eeg_harmonization.tex

\documentclass[preprint,12pt]{elsarticle}

%% For final journal submission use:
%% \documentclass[final,3p,times]{elsarticle}

\usepackage{amsmath}
\usepackage{amssymb}
\usepackage{graphicx}
\usepackage{booktabs}
\usepackage{multirow}
\usepackage{hyperref}
\usepackage{lineno}
\usepackage{microtype}
\usepackage{xcolor}
\usepackage{subcaption}

%% Uncomment for review mode:
% \linenumbers

\journal{Biomedical Signal Processing and Control}

% -----------------------------------------------------------------------
\begin{document}

\begin{frontmatter}

\title{(V$_{04}$) Harmonisation of Multi-Site EEG Data for Clinical Neuroscreening:
       Quantifying the Site Effect and Comparing Correction Strategies}

\author[aff1]{Kacper Chorzela}
\author[aff1]{Julia Pelczar}
\author[aff1]{Jaros\l{}aw {\.Z}ygierewicz\fnref{fn1}\corref{cor1}}
\ead{jarekz@fuw.edu.pl}

\cortext[cor1]{Corresponding author.}
\fntext[fn1]{ORCID: 0000-0002-7536-0735}

\affiliation[aff1]{organization={Faculty of Physics, University of Warsaw},
             addressline={Pasteura 5},
             city={Warsaw},
             postcode={02-093},
             country={Poland}}

% -----------------------------------------------------------------------
\begin{abstract}
Pooling electroencephalographic (EEG) recordings from multiple clinical
centres introduces systematic, non-biological variance---commonly
termed the \emph{site effect} or \emph{batch effect}---that can
severely degrade the generalisation of machine-learning diagnostic
models.
We investigate the site effect in a large-scale
collection of 54\,779 clinical EEG recordings from 30 Polish hospitals,
and evaluate a range of harmonisation strategies spanning simple
site-wise standardisation, the ComBat and CovBat algorithms, and
PCA-based feature-space pre-orthogonalisation applied as a
preprocessing step before any marginal correction method.
A CatBoost gradient-boosted classifier trained to identify the
recording institution from a 2\,850-dimensional feature set
(Riemannian covariance projections, multitaper power spectra, and
inter-electrode coherence estimates) achieved an overall
Matthews Correlation Coefficient (MCC) of $\mathbf{[TO\ BE\ FILLED]}
\pm \mathbf{[TO\ BE\ FILLED]}$, confirming that site identity is
strongly encoded in the feature space.
SHAP explainability analysis identified inter-electrode coherence
features---in particular occipital O1--O2 coherence and beta-band
temporal coherence---as the primary carriers of site-discriminative
information, implicating differences in recording protocols and
patient-preparation procedures across institutions.
Applying marginal correction methods (site-wise standardisation,
neuroCombat, and CovBat) without pre-orthogonalisation consistently
\emph{increased} the site-classification MCC towards a near-ceiling
value, demonstrating a harmonisation paradox: univariate feature
correction reorganises rather than eliminates the multivariate site
structure encoded in feature covariances.
Notably, logistic regression applied to the same harmonised features
showed \emph{reduced} site discriminability, demonstrating that the
paradox is classifier-dependent: it arises from the interaction
between residual multivariate covariance structure and nonlinear
classifier expressiveness.
A sensitivity analysis across PCA retained-variance levels (100\%,
99\%, 95\%, 90\%, 80\%) revealed that retaining 90\% of variance
yielded the best trade-off between site-invariance and clinical
utility.
PCA pre-orthogonalisation, in combination with any of the marginal
correction methods, yields more uniform cross-site neuroscreening
performance than marginal correction alone.
Despite the paradox, all harmonisation strategies produced modest but
consistent improvements in the downstream neuroscreening task under
both 5-fold cross-validation and a rigorous Leave-One-Site-Out (LOSO)
protocol.
These results indicate that the site effect constitutes a complex,
high-dimensional challenge that standard marginal correction methods
cannot fully resolve, and motivate signal-level adversarial
harmonisation as a complementary direction.
\textbf{[DANN RESULTS --- Julia to fill]}
\end{abstract}

\begin{keyword}
EEG \sep site effect \sep batch effect \sep data harmonisation \sep
ComBat \sep CovBat \sep PCA pre-orthogonalisation \sep neuroscreening
\sep multi-centre study \sep machine learning \sep SHAP \sep DANN
\end{keyword}

\end{frontmatter}

% -----------------------------------------------------------------------
\section{Introduction}
\label{sec:introduction}

Electroencephalography (EEG) is a non-invasive, cost-effective
neurophysiological technique that holds substantial promise for
scalable, data-driven clinical diagnostics~\cite{poziomska2026quantity}.
The increasing availability of large, routinely collected clinical EEG
archives has opened new avenues for training machine-learning (ML)
classifiers capable of detecting neurological abnormalities at a
population scale~\cite{mushtaq2025eeg}.
A prerequisite for such systems to be deployable in practice, however,
is the ability to generalise across acquisition environments: a model
trained on recordings from one set of hospitals must perform reliably
on data from institutions not represented during training.

This requirement is fundamentally challenged by the \emph{site
effect}---systematic, non-biological variance arising from
inter-centre differences in EEG hardware, electrode placement
conventions, amplifier settings, recording protocols, patient-
preparation procedures, and patient
demography~\cite{jaramillo2024combat,li2022harmnqeeg}.
When ML models are trained on multi-centre data without addressing
the site effect, they can exploit institution-specific signatures in
the feature space rather than learning clinically discriminative
patterns, leading to marked degradation of performance when the model
encounters data from a previously unseen
hospital~\cite{kiessner2023extended,mellot2023meeg}.
As large, heterogeneous clinical EEG corpora become the natural
substrate for the next generation of diagnostic algorithms, the site
effect has emerged as one of the principal obstacles to building
generalisable neuroscreening models~\cite{dong2025normative,
mushtaq2025eeg}.

Statistical harmonisation methods, most prominently the ComBat
family~\cite{johnson2007adjusting,fortin2017harmonization,
fortin2018cortical}, were developed to address the analogous
batch-effect problem in genomics and neuroimaging, and have more
recently been applied to EEG
features~\cite{jaramillo2024combat,li2022harmnqeeg}.
These methods model the site effect as additive and multiplicative
shifts in each feature's marginal distribution and remove them via
empirical-Bayes estimation.
CovBat~\cite{hu2023image} extends this framework to the multivariate
setting by additionally harmonising batch-specific structure in the
principal-component space of the residual covariance matrix.
Although computationally lightweight and widely available, all such
methods share a fundamental constraint: they operate feature by
feature and may therefore leave intact site-related structure encoded
in inter-feature covariances---a particularly important limitation in
high-dimensional EEG feature spaces where inter-electrode
relationships carry substantial clinical
information~\cite{hu2023image,mellot2023meeg}.

An alternative paradigm embeds harmonisation directly into the
representation-learning objective by training a neural network to
produce features that are simultaneously predictive of the clinical
label and invariant to the recording
site~\cite{ganin2016domain,ozdenizci2020invariant}.
Domain Adversarial Neural Networks (DANNs)~\cite{ganin2016domain},
realised via a gradient reversal layer, have demonstrated improved
cross-subject generalisation on EEG classification tasks and are a
natural candidate for the multi-site clinical harmonisation
setting~\cite{bethge2022multidomain,liu2023multisource}.

The present study pursues two overarching goals.
First, we provide a rigorous characterisation of the site effect in
the ELM19 corpus of 54\,779 EEG recordings from 30 Polish hospitals
provided by Elmiko Biosignals---using gradient-boosted classification,
SHAP feature-importance analysis, and Cohen's~$d$ effect-size
measurements.
Second, we evaluate a range of harmonisation strategies, spanning
simple site-wise standardisation, neuroCombat, CovBat, and PCA
pre-orthogonalisation applied as a preprocessing step before marginal
correction, and critically assess the conditions under which these
methods succeed or fail.
Signal-level adversarial harmonisation via a DANN architecture
integrated with a multiple-instance learning (MIL) model is
introduced as a complementary direction and is evaluated in parallel.
\textbf{[DANN --- Julia to fill]}

The main contributions of this paper are as follows:
\begin{enumerate}
  \item We demonstrate that the site effect in a large clinical EEG
        corpus is strong and broadly distributed across all feature
        types, with a gradient-boosted classifier achieving
        $\mathrm{MCC} = \mathbf{[TO\ BE\ FILLED]}$ in a site
        identification task.
  \item We identify, via SHAP analysis and ablation experiments,
        that inter-electrode coherence features are the primary
        carriers of site-specific information, and discuss the likely
        procedural origins of this pattern.
  \item We reveal and characterise a \emph{harmonisation paradox}:
        marginal correction methods (site-wise standardisation,
        neuroCombat, CovBat) reduce univariate inter-site feature
        differences (Cohen's~$d$) yet simultaneously increase the
        site-classification MCC towards a near-ceiling value,
        demonstrating that marginal correction reorganises rather
        than eliminates the multivariate site structure.
  \item We show that PCA pre-orthogonalisation of the feature space
        before marginal correction substantially attenuates the
        harmonisation paradox and stabilises cross-site neuroscreening
        performance across institutions.
  \item We demonstrate that the harmonisation paradox is
        classifier-dependent: a linear classifier (logistic regression)
        does not exhibit the paradox, while a nonlinear gradient-boosted
        classifier does, establishing that the paradox arises from
        classifier capacity to exploit residual covariance structure.
  \item We validate that the site effect is robust across five
        preprocessing filter configurations, confirming it reflects
        genuine inter-site differences rather than preprocessing
        artefacts.
  \item \textbf{[DANN contribution --- Julia to fill]}
  \item We evaluate all conditions on a downstream neuroscreening
        task under both standard cross-validation and a rigorous
        LOSO protocol, and compare feature-level methods against a
        signal-level adversarial baseline.
\end{enumerate}

% -----------------------------------------------------------------------
\section{Related Work}
\label{sec:related}

\subsection{The Site Effect in Multi-Centre EEG Studies}
\label{subsec:site_effect}

Systematic, non-biological variability arising from differences in
data acquisition conditions is a well-established problem across
quantitative biomedical sciences, where it is collectively termed the
\emph{batch effect}~\cite{johnson2007adjusting}.
In EEG research, the equivalent phenomenon---reproducible,
institution-specific signatures reflecting measurement context rather
than neurological state---is termed the \emph{site
effect}~\cite{li2022harmnqeeg}.
As large multi-centre EEG corpora become increasingly central to
data-driven clinical neuroscience~\cite{dong2025normative,
mushtaq2025eeg}, the site effect has emerged as a primary obstacle to
building classifiers that generalise beyond the hospitals represented
in training data.

Site effects arise from a heterogeneous combination of technical and
procedural factors.
At the hardware level, differences in amplifier design, electrode
materials, impedance characteristics, and recording-pipeline transfer
functions influence the spectral and spatial properties of the
acquired signal~\cite{li2022harmnqeeg}.
Recording protocols introduce further variability: institutions differ
in the use of activation procedures, resting-state recording duration,
and conventions regarding eye-open versus eye-closed
epochs~\cite{jaramillo2024combat}.
Patient-handling practices---scalp preparation, conductive gel
application, patient instruction, and alertness
management---constitute an additional source of inter-site divergence
that is rarely documented but may substantially affect features
sensitive to muscle artefact, slow drift, and vigilance
state~\cite{jaramillo2024combat,mushtaq2025eeg}.
Demographic and clinical case-mix differences across institutions
further contribute to apparent site differences even when acquisition
protocols are nominally identical~\cite{dong2025normative}.
Because these sources act simultaneously and interact in complex ways,
the site effect is not a simple additive shift in individual features
but a high-dimensional, multivariate structure distributed across the
full feature space~\cite{mellot2023meeg}.

The degree to which site identity is encoded in EEG feature spaces
has been documented quantitatively.
Li~et~al.~\cite{li2022harmnqeeg} demonstrated systematic batch
effects across 12 devices and 9 countries in the HarMNqEEG corpus of
1\,564 subjects, showing that cross-spectral features cluster by
acquisition device before harmonisation.
Jaramillo-Jimenez~et~al.~\cite{jaramillo2024combat} confirmed, using
$t$-SNE visualisation and Kruskal--Wallis tests, that spectral
features from five independent resting-state EEG datasets cluster by
site, with particularly prominent effects in the aperiodic offset and
oscillatory alpha power.
In the context of clinical EEG pathology decoding,
Kiessner~et~al.~\cite{kiessner2023extended} observed marked
performance degradation when a deep neural network trained on one
clinical corpus was evaluated on data from a different centre.
Mellot~et~al.~\cite{mellot2023meeg} showed that even
covariance-based EEG representations suffer substantial dataset shift
across acquisition devices and populations.

\subsection{Statistical Harmonisation Methods}
\label{subsec:statistical_harmonisation}

The most widely adopted family of harmonisation methods operates as
post-hoc corrections applied to pre-extracted feature representations.
Their appeal lies in computational simplicity, interpretability, and
the availability of well-maintained open-source implementations.
Their fundamental constraint is that they operate feature by feature
and therefore cannot capture higher-order, multivariate inter-site
structure~\cite{hu2023image,jaramillo2024combat}.

The canonical method is ComBat, originally proposed by
Johnson~et~al.~\cite{johnson2007adjusting} for batch-effect
correction in microarray gene-expression data.
ComBat models the observed feature value as the sum of a global
mean, a biological covariate effect, and additive and multiplicative
site-specific shifts, with site parameters estimated via empirical
Bayes to stabilise estimates at small
sites~\cite{johnson2007adjusting}.
Harmonised values are obtained by subtracting the estimated site
contributions while restoring the biological covariate effects.
Fortin~et~al.~\cite{fortin2017harmonization,fortin2018cortical}
demonstrated that this framework effectively removes scanner effects
from diffusion tensor imaging features and cortical thickness
measurements across multiple scanners, establishing ComBat as the
standard for retrospective neuroimaging harmonisation.

A growing ecosystem of ComBat extensions addresses specific
limitations of the original model.
ComBat-GAM (neuroHarmonize~\cite{pomponio2020harmonization}) replaces
the linear covariate model with a generalised additive model to
handle non-linear biological covariate effects.
CovBat~\cite{hu2023image} represents a qualitatively distinct
extension: it first applies standard ComBat to correct marginal mean
and variance shifts, and then additionally harmonises the
principal-component space of the residual covariance matrix, thereby
addressing batch effects that reside in the inter-feature covariance
structure rather than in individual feature distributions.
CovBat has been shown to outperform standard ComBat in both
multivariate statistical evaluations and downstream machine-learning
metrics in neuroimaging data~\cite{hu2023image}.
A comprehensive taxonomy and benchmark of ComBat variants is provided
by Hu~et~al.~\cite{hu2023image}.

In the EEG domain, Jaramillo-Jimenez~et~al.~\cite{jaramillo2024combat}
evaluated four ComBat variants on pooled resting-state EEG data from
five independent datasets and found that all variants reduced
batch-related dispersion in spectral features while preserving
age-related associations.
Li~et~al.~\cite{li2022harmnqeeg} demonstrated that ComBat-based
normalisation substantially reduced device-related clustering across
12 devices and 9 countries, yielding the HarMNqEEG normative database.

Despite widespread adoption, ComBat has well-characterised
limitations.
Most critically, adjusting mean and variance independently for each
feature leaves intact any site-related structure residing in the
\emph{covariance} between features~\cite{hu2023image}.
CovBat was designed specifically to address this limitation, but
its principal-component decomposition step may be unstable when
individual sites have small sample sizes~\cite{hu2023image}.
A second limitation shared by all ComBat-family methods is that
biological covariates must be available and correctly specified at
harmonisation time; misspecification can inadvertently remove genuine
biological variation alongside the site
effect~\cite{hu2023image,fortin2017harmonization}.
Additionally, ComBat requires all sites to be present at
parameter-estimation time and cannot directly accommodate recordings
from a site not represented in the training data.

Normative modelling offers a conceptually distinct approach to
multi-site heterogeneity~\cite{dong2025normative}: rather than
correcting features toward a common distribution, normative models
estimate reference trajectories from a large healthy population and
express each individual's measurements as a deviation from the
age-adjusted norm.
Dong~et~al.~\cite{dong2025normative} demonstrated this framework for
large-scale multisite EEG data, using ComBat to remove residual batch
effects prior to constructing normative models; the resulting
quantile-rank features significantly outperformed raw features in
predicting attention-related outcomes.
Normative modelling nevertheless relies on the same underlying ComBat
correction as a preprocessing step and therefore inherits its
limitations with respect to multivariate site structure.

\subsection{Deep Learning and Adversarial Harmonisation}
\label{subsec:dann_related}

An alternative paradigm embeds harmonisation directly into the
representation-learning objective, training a neural network to
produce features that are simultaneously predictive of the clinical
label and invariant to the recording site.
This belongs to the broader field of \emph{domain
adaptation}~\cite{ganin2016domain}, which seeks to train models that
generalise across distributional shifts between source and target
domains.

The canonical architecture is the Domain Adversarial Neural Network
(DANN)~\cite{ganin2016domain}, which consists of a shared feature
extractor, a label predictor, and a domain (site) discriminator.
The minimax training objective is realised by a Gradient Reversal
Layer (GRL), which multiplies the gradient from the discriminator by
$-\lambda$ during backpropagation, forcing the feature extractor to
produce representations from which site identity cannot be
recovered~\cite{ganin2016domain}.
Compared with ComBat-family methods, adversarial domain adaptation
offers structural advantages: harmonisation is task-aware and embedded
in representation learning rather than applied post-hoc; the domain
classifier operates on the full latent representation and can in
principle capture multivariate inter-site structure; and the framework
does not require all sites to be observed during parameter estimation.

Adversarial representation learning has been applied widely in the
EEG literature.
\"{O}zdenizci~et~al.~\cite{ozdenizci2020invariant} demonstrated
significant improvements in cross-subject generalisation on motor
imagery tasks by training a convolutional network with an adversarial
subject-domain classifier.
Liu~et~al.~\cite{liu2023multisource} proposed MTDANN, which trains
separate DANNs for each source--target pair using
EEGNet~\cite{lawhern2018eegnet} as the shared feature extractor and
combines predictions via a weighted fusion mechanism; MTDANN
outperformed single-source DANN and transfer learning baselines on
three public motor imagery datasets.
Bethge~et~al.~\cite{bethge2022multidomain} showed that a unified
adversarial invariance objective across multiple emotion recognition
datasets suppressed data-source-related information leakage by 35\%
while maintaining classification performance, demonstrating that the
DANN principle extends naturally to the multi-dataset and multi-site
setting.
The adversarial approach we develop in the present work
(Section~\ref{subsec:approach2_dann}) integrates the DANN objective
with the MIL architecture of Poziomska~et~al.\
\cite{poziomska2026quantity}, which was designed specifically for
multi-recording EEG pathology decoding at scale.

% -----------------------------------------------------------------------
\section{Materials and Methods}
\label{sec:methods}

\subsection{Dataset}
\label{subsec:dataset}

\subsubsection{The ELM19 Corpus}

We use the ELM19 database~\cite{poziomska2026quantity}, a large-scale
multi-centre EEG corpus provided by Elmiko Biosignals sp.~z~o.o.\
(Milanówek, Poland).
ELM19 comprises 55\,787 clinical EEG recordings collected from
39 medical institutions across Poland during routine clinical
diagnostics.
The dataset spans a broad spectrum of patient cases, from standard
neurological examinations to assessments for diverse clinical
conditions.
Each recording is annotated with a binary normal/pathological label
derived from the accompanying EEG report following the procedure
detailed in~\cite{poziomska2026quantity}.
Patient ages range from 16 to 70 years; 52.9\% of recordings are from
female patients; and 52.2\% of recordings are labelled as
pathological.
EEG signals were acquired using the international 10--20 electrode
placement system with 19 channels (Fp1, Fp2, F7, F3, Fz, F4, F8,
T3, C3, Cz, C4, T4, T5, P3, Pz, P4, T6, O1, O2), with sampling
rates varying between 200 and 500\,Hz across institutions.

\subsubsection{Site Selection and the
  \texorpdfstring{$\mathrm{ELM}_s$}{ELMs} Subset}

Sites contributing fewer than 50 normal recordings were excluded to
ensure sufficient within-site representation for harmonisation
parameter estimation; such sites were deemed insufficient for reliable
characterisation.
The threshold was determined by inspecting the distribution of normal
recordings across all institutions
(Figure~\ref{fig:institution_histograms}).
After this filtering step, 9 institutions were excluded, and the
resulting dataset---hereafter $\mathrm{ELM}_s$---consists of 54\,779
EEG recordings from 30 medical institutions, comprising 26\,568
normal (48.5\%) and 28\,211 pathological (51.5\%) recordings from
28\,960 female subjects (52.9\%).

\begin{figure}[htbp]
  \centering
  \begin{subfigure}[t]{0.48\textwidth}
    \centering
    \includegraphics[width=\textwidth]{%
      Figs/ELM19_excluded_institutions_normal_recordings.png}
    \caption{Number of recordings classified as normal per institution.}
    \label{fig:normal_hist_excluded}
  \end{subfigure}
  \hfill
  \begin{subfigure}[t]{0.48\textwidth}
    \centering
    \includegraphics[width=\textwidth]{%
      Figs/ELM19_excluded_institutions_all_recordings.png}
    \caption{Total number of recordings (normal and pathological)
      per institution.}
    \label{fig:total_hist_excluded}
  \end{subfigure}
  \caption{Distribution of recordings across the 39 medical
    institutions in the ELM19 corpus. Institutions below the
    50-normal-recording threshold (dashed line) were excluded from
    further analysis.}
  \label{fig:institution_histograms}
\end{figure}

\subsubsection{Analysis Subset: Normal Recordings
  (\texorpdfstring{$\mathrm{ELM}_n$}{ELMn})}

To isolate the site effect from variance attributable to clinical
pathology, all primary analyses---exploratory data analysis, site
classification, harmonisation evaluation, and feature importance---were
conducted on the subset of recordings labelled as normal, hereafter
referred to as $\mathrm{ELM}_n$.
Restricting to normal recordings eliminates the confounding influence
of pathology-related signal variability, enabling a cleaner
characterisation of between-site baseline differences.
\subsubsection{Acquisition Metadata}
\label{subsubsec:acquisition_metadata}

The 30 institutions in $\mathrm{ELM}_s$ used two Elmiko hardware
models (1032 and 1042) and a range of software versions, with some
institutions changing software versions during the data collection
period (Table~\ref{tab:acquisition_metadata}).
Sampling frequencies were predominantly $\sim$250\,Hz, with isolated
600\,Hz anomalies at certain sites.
Acquisition-level filter settings---including highpass and lowpass
cutoff frequencies---also varied across institutions.
\textbf{[TO BE FILLED: brief summary of hardware/software
distribution across sites.]}
These metadata provide concrete evidence that the site effect has
identifiable technical origins beyond purely demographic or clinical
differences.
A subset of the institutions are specialist hospitals (e.g.\
children's hospitals), introducing additional patient population
heterogeneity.

\begin{table}[htbp]
  \centering
  \caption{Acquisition metadata for the 30 institutions in
    $\mathrm{ELM}_s$. Hardware model, software version, and
    specialisation status are shown where available.}
  \label{tab:acquisition_metadata}
  \begin{tabular}{lcccc}
    \toprule
    \textbf{Institution} & \textbf{HW Model} &
    \textbf{SW Version} & \textbf{Children's} &
    \textbf{Specialisation} \\
    \midrule
    \multicolumn{5}{c}{\textbf{[TO BE FILLED from
      Lista\_baz\_informacje.xlsx]}} \\
    \bottomrule
  \end{tabular}
\end{table}

The full $\mathrm{ELM}_s$ dataset was used only in the downstream
neuroscreening experiment.

\subsection{EEG Preprocessing}
\label{subsec:preprocessing}

Raw EEG signals were preprocessed using
MNE-Python~\cite{gramfort2013mne} following the pipeline described
in~\cite{poziomska2026quantity}.
A Butterworth high-pass filter (passband edge 0.5\,Hz, stopband edge
0.1\,Hz) was applied to remove slow drift, together with a low-pass
filter (cutoff 40\,Hz) to suppress high-frequency muscle artefacts.
A notch filter ($f_0 = 50$\,Hz, quality factor $Q = 5$) attenuated
power-line interference; for institutions located in the United
States, a 60\,Hz notch filter was applied instead.
All signals were re-referenced to a common average reference and
resampled to 100\,Hz to ensure a uniform sampling rate across all
institutions.
Continuous recordings were segmented into non-overlapping 6-second
frames (600 samples at 100\,Hz).
Frames containing channels with median absolute deviation below
$1.0\,\mu\text{V}$ (flat-line detection) or peak-to-peak voltages
exceeding $800\,\mu\text{V}$ were discarded as artefacts.
Per-institution frame rejection rates were computed to characterise
signal quality heterogeneity across sites
(Section~\ref{subsec:exploratory}).

\subsection{Feature Extraction}
\label{subsec:features}

Three groups of features were extracted from each 6-second frame,
following the methodology of~\cite{poziomska2026quantity}.
Power and coherence features were aggregated per recording using the
median across all valid frames.
Covariance matrices were aggregated using the Riemannian geometric
mean (see below).

\subsubsection{Riemannian Covariance Features}

For each frame, the $19 \times 19$ channel covariance matrix was
computed and frames within a recording were averaged using the
Riemannian geometric mean via the \texttt{pyRiemann} library.
Because covariance matrices lie on a Riemannian manifold, the average
matrix was projected onto its tangent space to obtain a Euclidean
representation; the lower triangular elements were extracted,
yielding \textbf{190 features} per recording that capture the spatial
covariance structure of the EEG.

\subsubsection{Multitaper Power Features}

Power spectral densities (PSD) were estimated for each channel using
the multitaper method~\cite{thomson1982spectrum} and integrated across
14 custom overlapping frequency bands covering 0.5--40\,Hz:
$[0.5, 2]$, $[1, 3]$, $[2, 4]$, $[3, 6]$, $[4, 8]$, $[6, 10]$,
$[8, 13]$, $[10, 15]$, $[13, 18]$, $[15, 21]$, $[18, 24]$,
$[21, 27]$, $[24, 30]$, and $[27, 40]$\,Hz, spanning the standard
delta, theta, alpha, and beta EEG bands.
PSDs were normalised such that the sum across all channels and bands
equalled one within each frame, yielding \textbf{266 power features}
per recording.

\subsubsection{Inter-Electrode Coherence Features}

Band-wise magnitude-squared coherence was computed for all
$\binom{19}{2} = 171$ inter-electrode pairs in each of the 14
frequency bands:
\begin{equation}
  C_{xy}(f_b) = \frac{\bigl|S_{xy}(f_b)\bigr|^2}
                     {S_{xx}(f_b)\cdot S_{yy}(f_b)},
  \label{eq:coherence}
\end{equation}
where $S_{xy}(f_b)$ is the cross-spectral density and $S_{xx}(f_b)$,
$S_{yy}(f_b)$ are the auto-spectral densities of channels $x$ and $y$
in band $f_b$, yielding $171 \times 14 = \mathbf{2\,394}$
\textbf{coherence features} per recording.

\subsubsection{Final Feature Vector}

Concatenating the three groups yields a final feature vector of
$190 + 266 + 2\,394 = \mathbf{2\,850}$ features per recording,
capturing spatial covariance in the time domain, spectral power
across a wide frequency range, and functional connectivity between
all electrode pairs.

\subsection{Approach~1: Feature-Level Harmonisation}
\label{subsec:approach1}

All feature-level methods were applied independently to each of the
2\,850 features.
Harmonisation parameters were estimated exclusively on the training
portion of each data split and subsequently applied to both training
and test data to prevent data leakage.

\subsubsection{Site-Wise Z-Score Standardisation}
\label{subsubsec:sitewise}

As a conceptually simple baseline, each feature was standardised to
zero mean and unit variance using within-site statistics estimated
from the training split.
For subject $i$ from site $s(i)$ and feature $j$:
\begin{equation}
  z_{ij} = \frac{x_{ij} - \hat{\mu}_{s(i),j}}
                {\hat{\sigma}_{s(i),j}},
  \label{eq:sitewise}
\end{equation}
where $\hat{\mu}_{s(i),j}$ and $\hat{\sigma}_{s(i),j}$ are the
within-site mean and standard deviation of feature $j$ estimated on
the training data.
This transformation equalises the first two moments of each feature
distribution across sites but makes no assumptions about the
structure of the site effect, does not account for biological
covariates, and leaves inter-feature covariance structure entirely
uncorrected.

\subsubsection{neuroCombat}
\label{subsubsec:neurocombat}

The neuroCombat implementation of the ComBat
algorithm~\cite{johnson2007adjusting,fortin2017harmonization} was
applied with the hospital identifier as the batch variable and patient
age and sex as biological covariates.
ComBat models the observed feature value $y_{ij}$ for subject $i$ and
feature $j$ as
\begin{equation}
  y_{ij} = \alpha_j + \mathbf{x}_i^{\top}\boldsymbol{\beta}_j
           + \gamma_{s(i),j} + \delta_{s(i),j}\,\varepsilon_{ij},
  \label{eq:combat}
\end{equation}
where $\alpha_j$ is the global intercept, $\mathbf{x}_i$ is the
vector of biological covariates with coefficients
$\boldsymbol{\beta}_j$, $\gamma_{s(i),j}$ is an additive site
effect, $\delta_{s(i),j}$ is a multiplicative site effect, and
$\varepsilon_{ij}$ is residual noise.
Site parameters are estimated via empirical Bayes, pooling information
across features to stabilise estimates at small
sites~\cite{johnson2007adjusting}.
Harmonised values are recovered by subtracting the estimated site
contributions while restoring the biological covariate effects:
\begin{equation}
  y^*_{ij} = \frac{y_{ij} - \hat{\alpha}_j
             - \mathbf{x}_i^{\top}\hat{\boldsymbol{\beta}}_j
             - \hat{\gamma}_{s(i),j}}{\hat{\delta}_{s(i),j}}
             + \hat{\alpha}_j
             + \mathbf{x}_i^{\top}\hat{\boldsymbol{\beta}}_j.
  \label{eq:combat_corrected}
\end{equation}
Like site-wise standardisation, neuroCombat corrects each feature
independently and therefore leaves inter-feature covariance structure
unaffected.
We also evaluated standard ComBat without biological covariates;
because the results were numerically indistinguishable from
neuroCombat with age and sex covariates in all experiments, only the
neuroCombat variant is reported hereafter.

\subsubsection{CovBat}
\label{subsubsec:covbat}

CovBat~\cite{hu2023image} extends neuroCombat to address site-related
structure residing in the inter-feature covariance.
The method first applies standard ComBat (Equations~\ref{eq:combat}
and~\ref{eq:combat_corrected}) to correct marginal mean and variance
shifts, and then performs a second harmonisation step in the
principal-component space of the residual covariance matrix.
Specifically, the ComBat-corrected residuals are projected onto their
leading principal components; site-specific shifts in this PC space
are estimated and removed, and the corrected PCs are projected back
to the original feature space.
This two-stage procedure targets both marginal and covariance-level
site effects, making CovBat the most comprehensive of the
feature-level methods evaluated here.
We used the implementation of Hu~et~al.~\cite{hu2023image} with
patient age and sex as biological covariates.

\subsubsection{PCA Pre-Orthogonalisation}
\label{subsubsec:pca}

As an additional preprocessing step---applied \emph{before} any of
the marginal correction methods above---we evaluated PCA
pre-orthogonalisation of the feature space.
PCA projections were computed on the training data at six
retained-variance levels: 100\% (all components), 99\%, 95\%, 90\%,
and 80\%, and applied to both training and test data prior to
harmonisation.
The full-variance projection serves as a pure decorrelation step;
reduced-variance projections additionally remove low-variance
components that may encode noise or site-specific structure.
Because the resulting PC features are mutually uncorrelated by
construction, PCA pre-orthogonalisation transforms the feature space
into one in which inter-feature covariance is identically zero.
Marginal correction methods subsequently applied to this transformed
space therefore operate in a setting where the inter-feature
covariance structure---identified in Section~\ref{subsubsec:paradox}
as the mechanism underlying the harmonisation paradox---is
destroyed.
PCA pre-orthogonalisation was evaluated both as a standalone
condition (without any subsequent marginal correction) and in
combination with site-wise standardisation, neuroCombat, and CovBat.
Results across all variance levels are reported in
Section~\ref{subsec:pca_sweep}.

\subsubsection{Downstream Classifier}
\label{subsubsec:classifier}

Original or harmonised feature vectors were passed to a
gradient-boosted ensemble (GBE) model composed of 30 individual
CatBoost~\cite{prokhorenkova2018catboost} classifiers, identical to
the architecture described in~\cite{poziomska2026quantity}.
CatBoost was selected for its effectiveness on high-dimensional
tabular data, robustness to overfitting, and native support for
multi-class classification.
Hyperparameters were tuned via Bayesian optimisation using the
Optuna framework on the training fold with respect to AUC.
[TO BE FILLED: final hyperparameter values.]

As a linear baseline, $\ell_2$-regularised logistic regression was
applied to the same feature vectors.
Comparing a linear classifier that cannot exploit inter-feature
covariance structure against a nonlinear gradient-boosted ensemble
isolates the role of classifier expressiveness in the harmonisation
paradox.

% TODO(Julia): Review and update this section with actual implementation details
\subsection{Approach~2: Signal-Level Adversarial Harmonisation}
\label{subsec:approach2_dann}

Feature-level harmonisation methods operate on pre-extracted,
hand-crafted representations and therefore cannot address site-related
structure residing in the inter-feature covariance, as demonstrated
empirically by the harmonisation paradox in
Section~\ref{subsubsec:paradox}.
To overcome this limitation, we develop a signal-level approach based
on Domain Adversarial Neural Networks~\cite{ganin2016domain}, which
learns site-invariant representations end-to-end from raw EEG frames.
The architecture follows the multiple-instance learning (MIL)
paradigm of Poziomska~et~al.~\cite{poziomska2026quantity} and consists
of three functional components.

\subsubsection{Network Architecture}
\label{subsubsec:dann_arch}

\paragraph{Frame encoder.}
An EEGNet-based~\cite{lawhern2018eegnet} convolutional feature
extractor $G_f$ maps each 6-second EEG frame
$\mathbf{x} \in \mathbb{R}^{C \times T}$ (19 channels, 600 time
points at 100\,Hz) to a frame-level latent representation
$\mathbf{f} \in \mathbb{R}^d$.
The encoder performs separate temporal and spatial convolutions and
is well-established as a compact, sample-efficient encoder for
multi-channel EEG~\cite{lawhern2018eegnet,poziomska2026quantity}.

\paragraph{Attention-based recording aggregation.}
Because clinical pathology may manifest in only a subset of the
frames within a recording, frame-level representations are aggregated
using the attention module of~\cite{poziomska2026quantity}, which
follows the non-gated attention formulation and incorporates a
query--key--value structure.
The module assigns a scalar attention weight $a_t$ to each frame and
computes the recording-level representation as
\begin{equation}
  \mathbf{h} = \sum_{t=1}^{T_r} a_t\, \mathbf{f}_t,
  \qquad
  a_t = \frac{\exp\!\bigl(\mathbf{w}^\top
              \tanh(\mathbf{V}\mathbf{f}_t)\bigr)}
             {\sum_{\tau=1}^{T_r}
              \exp\!\bigl(\mathbf{w}^\top
              \tanh(\mathbf{V}\mathbf{f}_\tau)\bigr)},
  \label{eq:attention}
\end{equation}
where $T_r$ is the number of valid frames in the recording, and
$\mathbf{V}$ and $\mathbf{w}$ are learned parameters.
Frames carrying stronger evidence of pathology are expected to receive
higher weights~\cite{poziomska2026quantity}.

\paragraph{Pathology classifier and site discriminator.}
The recording-level representation $\mathbf{h}$ is passed
simultaneously to a pathology classifier $G_y$---a two-layer fully
connected network producing a binary normal/pathological
prediction---and to a site discriminator $G_d$, a fully connected
network predicting membership among the 30 recording institutions.

\subsubsection{Training Objective}
\label{subsubsec:dann_training}

Training minimises the saddle-point objective
\begin{equation}
  \mathcal{L}(\theta_f,\theta_a,\theta_y,\theta_d) =
      \mathcal{L}_y\!\bigl(G_y(\mathbf{h}),y\bigr)
      - \lambda\,\mathcal{L}_d\!\bigl(G_d(\mathbf{h}),d\bigr),
  \label{eq:dann_minet}
\end{equation}
where $\theta_f$, $\theta_a$, $\theta_y$, and $\theta_d$ are the
parameters of the frame encoder, attention module, classifier, and
discriminator, respectively; $\mathcal{L}_y$ and $\mathcal{L}_d$ are
cross-entropy losses for pathology classification and site
discrimination; and $\lambda$ controls the trade-off between domain
confusion and classification accuracy.
The minimax objective is implemented via a Gradient Reversal
Layer~\cite{ganin2016domain} inserted between the attention module
and the site discriminator.
The domain-adaptation weight $\lambda$ is annealed according to the
progressive schedule of Ganin~et~al.~\cite{ganin2016domain}:
\begin{equation}
  \lambda_p = \frac{2}{1 + \exp(-\gamma p)} - 1, \quad \gamma = 10,
  \label{eq:lambda_schedule}
\end{equation}
where $p \in [0,1]$ denotes the fraction of total training epochs
completed.
This schedule suppresses adversarial pressure in early training,
allowing the encoder and classifier to first develop discriminative
representations before the site-invariance constraint is gradually
enforced.

\subsubsection{Training Procedure}
\label{subsubsec:dann_procedure}

The encoder is pretrained on individual frames in a single-instance
fashion---identical to the siNet pretraining of~\cite{poziomska2026quantity}---before the full MIL network is trained
end-to-end with the adversarial objective.
Pretraining initialises the encoder with discriminative frame-level
features and consistently yields better performance than training
from scratch~\cite{poziomska2026quantity}.
We evaluate two model variants: MINet (pretraining only, no
adversarial training, $\lambda \equiv 0$), which serves as a
signal-level baseline, and MINet-DANN (full adversarial training).
For the held-out site in the LOSO protocol, the calibration sample
of $N = 30$ normal recordings is used during both pretraining and
adversarial fine-tuning to expose the harmonisation objective to
the target site's distribution.

\subsection{Evaluation Protocol}
\label{subsec:evaluation}

\subsubsection{Quantifying the Site Effect}
\label{subsubsec:site_quant}

The site effect was quantified using a supervised site-classification
task: a CatBoost classifier was trained to predict the recording
institution from the 2\,850-dimensional feature vector on
$\mathrm{ELM}_n$, using 5-fold stratified cross-validation with
stratification by hospital ID, and, as a linear control, an
$\ell_2$-regularised logistic regression classifier was trained on
the same features.
The Matthews Correlation Coefficient (MCC) was used as the primary
metric because it provides a balanced measure of multi-class
classification quality under class imbalance.
For a $K$-class problem with confusion matrix
$\mathbf{C} \in \mathbb{N}^{K \times K}$, the multiclass MCC is:
\begin{equation}
  \mathrm{MCC} = \frac{c\,s - \sum_{k=1}^{K} p_k t_k}
    {\sqrt{\!\left(s^2 - \sum_{k=1}^{K} p_k^2\right)
           \!\left(s^2 - \sum_{k=1}^{K} t_k^2\right)}},
  \label{eq:mcc}
\end{equation}
where $c = \sum_k C_{kk}$, $s = \sum_{i,j} C_{ij}$,
$t_k = \sum_j C_{kj}$, and $p_k = \sum_i C_{ik}$.
Per-hospital MCC values were computed using a one-vs-rest approach.
The same site-classification task was repeated after applying each
harmonisation condition to assess how correction affected multivariate
site discriminability.
% TODO(Julia): Verify this matches your actual evaluation protocol
\subsubsection*{Site Discriminability of Adversarial Representations}
\label{subsubsec:dann_site_probe}

To enable a direct comparison with the feature-level site-classification
described above, we assess the residual
site information retained in the latent representations produced by the
MINet and MINet-DANN models.
After training on $\mathrm{ELM}_s$, each model is used in inference mode
to compute the recording-level representation $\mathbf{h}$ for every
recording in $\mathrm{ELM}_n$.
A linear probe---a logistic regression classifier with $\ell_2$
regularisation---is then trained on these representations to predict
site identity, using the same 5-fold stratified cross-validation
protocol applied to the feature-level methods.
The resulting MCC measures how much site identity is \emph{linearly
accessible} in the learned representation, evaluated on normal
recordings only, and is therefore directly comparable to the
site-classification MCC reported for all Approach~1 conditions in
Table~\ref{tab:site_class_harmonisation}.
Comparing MINet (no adversarial training) against MINet-DANN quantifies
the degree to which the adversarial objective suppresses site information
at the representation level, and comparing both against the
feature-level conditions establishes whether end-to-end representation
learning achieves a qualitatively different level of site invariance
than post-hoc marginal correction.

\subsubsection{Feature Importance Analysis}
\label{subsubsec:shap}

To identify which features were most responsible for site
discriminability, SHAP (SHapley Additive exPlanations) values were
computed for the trained CatBoost site-classification model using
the SHAP TreeExplainer:
\begin{equation}
  g(\mathbf{z}') = \phi_0 + \sum_{i=1}^{M} \phi_i z'_i,
  \label{eq:shap}
\end{equation}
where $\phi_0$ is the base value and $\phi_i$ is the SHAP value for
feature $i$.
Global feature importance was summarised by the mean absolute SHAP
value across all predictions and classes.
An ablation study trained models on each feature group individually
(covariance, power, coherence) and with each group sequentially
removed, to assess the independent and complementary contributions of
each feature type to site discriminability.

\subsubsection{Univariate Harmonisation Effectiveness: Cohen's~$d$}
\label{subsubsec:cohend}

The impact of harmonisation on individual feature distributions was
quantified using the standardised mean difference between institutions:
\begin{equation}
  d = \frac{\mu_1 - \mu_2}{\sigma_{\mathrm{pooled}}},
  \qquad
  \sigma_{\mathrm{pooled}} =
    \sqrt{\frac{(n_1-1)s_1^2 + (n_2-1)s_2^2}{n_1+n_2-2}},
  \label{eq:cohend}
\end{equation}
computed for each feature over every pair of institutions and
averaged across features to produce an all-vs-all pairwise summary.
Together with the multivariate site-classification MCC, Cohen's~$d$
provides a complementary univariate perspective on harmonisation
effectiveness.

\subsubsection{Clinical Neuroscreening Evaluation}
\label{subsubsec:clinical_eval}

Clinical utility was assessed on the binary pathology classification
task using the full $\mathrm{ELM}_s$ dataset with two protocols.

\paragraph{5-Fold Cross-Validation.}
The dataset was split into five folds stratified by hospital ID,
and performance (MCC and AUC) was averaged across folds.

\paragraph{Leave-One-Site-Out (LOSO).}
In each round, all recordings from one institution formed the test
set and the model was trained on the remaining 29 institutions.
Because harmonisation algorithms require training-data statistics
from the target site, a calibration sample of $N = 30$ normal
recordings from the held-out site was reserved prior to each round.
For harmonised conditions, this calibration sample was used jointly
with training data to fit the harmonisation parameters; for the
unharmonised condition, it was discarded to maintain identical train/
test splits across all conditions.
This protocol directly simulates deployment to a previously unseen
recording environment and is the most clinically relevant evaluation
setting in the present study.

\subsubsection{Robustness to Preprocessing Configuration}
\label{subsubsec:filter_robustness_eval}

To verify that the site effect is not an artefact of a specific
preprocessing pipeline, the site-classification experiment was
repeated under five filter configurations:
\textbf{[TO BE FILLED: list the 5 configs with brief descriptions of
what each changes]}.
Site-classification MCC was compared across configurations to assess
stability.

\subsubsection{PCA Variance Sensitivity Analysis}
\label{subsubsec:pca_sweep_eval}

To characterise the trade-off between site-invariance and clinical
utility as a function of retained variance, both site classification
and pathology classification were repeated at each PCA variance level
(100\%, 99\%, 95\%, 90\%, 80\%) combined with each harmonisation
method.

\subsubsection{Statistical Significance}
\label{subsubsec:statistical_tests}

Statistical significance of MCC differences across harmonisation
conditions was assessed using \textbf{[TO BE FILLED: Wilcoxon
signed-rank / paired $t$-test]} across the 30 LOSO folds.

% -----------------------------------------------------------------------
\section{Results}
\label{sec:results}

\subsection{Exploratory Analysis of Inter-Institutional Heterogeneity}
\label{subsec:exploratory}

Exploratory analysis of the $\mathrm{ELM}_n$ subset revealed
substantial heterogeneity across the 30 institutions in patient
demographics and recording protocols.
Age distributions varied markedly between sites: institutions such as
PIO, SL2, and TER\_L served predominantly young populations (median
ages 17.6, 19.5, and 17.8 years, respectively), while institutions
such as MOR and TOR served predominantly older patients (median ages
52.1 and 50.8 years, respectively).
The maximum difference in median age between any two institutions
exceeded 30 years (full statistics are given in
Table~\ref{tab:age_stats} in the Appendix).
Sex distributions were broadly balanced across most institutions,
with minor imbalances at a small number of sites.
Recording duration showed further heterogeneity: while the majority
of centres produced sessions of approximately 1\,000 seconds, notable
deviations were observed---institution KATMOJPRZ recorded sessions
lasting up to 8\,322 seconds, and ZOZLO showed a bimodal distribution
suggesting the use of multiple recording protocols within a single
institution.
This demographic and procedural heterogeneity motivates the use of
$\mathrm{ELM}_n$ for all site-effect analyses and provides a
qualitative indication that observed inter-site feature differences
reflect a mixture of acquisition and population-level sources.

The 30 institutions used two Elmiko hardware models (1032 and 1042)
and a range of software versions
(Table~\ref{tab:acquisition_metadata}).
\textbf{[TO BE FILLED: brief summary of HW/SW distribution.]}
Acquisition-level filter settings also varied:
\textbf{[TO BE FILLED: description of highpass/lowpass variation.]}
These metadata provide direct evidence that technical acquisition
differences contribute to the site effect.

Frame rejection rates during preprocessing varied substantially
across institutions, from \textbf{[TO BE FILLED]}\% to
\textbf{[TO BE FILLED]}\%, reflecting differences in recording
quality and electrode preparation practices across sites
(Figure~\ref{fig:rejection_rates}).

\begin{figure}[htbp]
  \centering
  % \includegraphics[width=\textwidth]{Figs/rejection_rates.pdf}
  \caption{Per-institution frame rejection rates during
    preprocessing, reflecting signal quality heterogeneity across
    the 30 institutions in $\mathrm{ELM}_n$.}
  \label{fig:rejection_rates}
\end{figure}

\subsection{Quantifying the Site Effect}
\label{subsec:site_classification}

\subsubsection{Site Classification Performance}

The CatBoost classifier trained on the $\mathrm{ELM}_n$ feature set
achieved an overall MCC of $\mathbf{[TO\ BE\ FILLED]} \pm
\mathbf{[TO\ BE\ FILLED]}$ across the 30 institutions in the 5-fold
stratified cross-validation, confirming that site identity is
strongly and reliably encoded in the feature space.
Per-institution one-vs-rest MCC scores are summarised in
Table~\ref{tab:per_site_mcc}; \textbf{[TO BE FILLED: number]} out of
30 institutions achieved average MCC scores above 0.7, with a small
subset of institutions (\textbf{[TO BE FILLED: list]}) being
considerably more difficult to classify, likely owing to distributional
similarity with neighbouring sites in the corpus.

\begin{table}[htbp]
  \centering
  \caption{Per-institution one-vs-rest MCC (mean $\pm$ SD) for the
    unharmonised site-classification task (5-fold CV on
    $\mathrm{ELM}_n$, all 2\,850 features).
    Institutions are ordered by decreasing raw-baseline MCC.}
  \label{tab:per_site_mcc}
  \begin{tabular}{lc@{\quad}lc}
    \toprule
    \textbf{Institution} & \textbf{MCC} &
    \textbf{Institution} & \textbf{MCC} \\
    \midrule
    MOR       & [TO BE FILLED] & PRZ      & [TO BE FILLED] \\
    B2K       & [TO BE FILLED] & KLU      & [TO BE FILLED] \\
    TOR       & [TO BE FILLED] & WLU      & [TO BE FILLED] \\
    KAL       & [TO BE FILLED] & MKW      & [TO BE FILLED] \\
    LUMICE    & [TO BE FILLED] & SRK      & [TO BE FILLED] \\
    KUD       & [TO BE FILLED] & OST      & [TO BE FILLED] \\
    KATMOJPRZ & [TO BE FILLED] & ZOZLO    & [TO BE FILLED] \\
    STG1      & [TO BE FILLED] & OTW      & [TO BE FILLED] \\
    Z04O      & [TO BE FILLED] & GAK      & [TO BE FILLED] \\
    SZC       & [TO BE FILLED] & TER\_L   & [TO BE FILLED] \\
    SLU       & [TO BE FILLED] & PIO      & [TO BE FILLED] \\
    PUS       & [TO BE FILLED] & AKS      & [TO BE FILLED] \\
    ARCHDAM   & [TO BE FILLED] & SL2      & [TO BE FILLED] \\
    LUX\_A    & [TO BE FILLED] & CHE      & [TO BE FILLED] \\
    overall   & [TO BE FILLED] & KIEG     & [TO BE FILLED] \\
              &                & CMD      & [TO BE FILLED] \\
    \bottomrule
  \end{tabular}
\end{table}

\subsubsection{Feature Importance: SHAP Analysis and Ablation}
\label{subsubsec:shap_results}

Table~\ref{tab:ablation} summarises site-classification MCC for models
trained on individual feature groups and with each group excluded.
\textbf{[TO BE FILLED: narrative on which group dominates and what
the exclusion experiments reveal.]}
The ablation results are corroborated by the SHAP analysis:
\textbf{[TO BE FILLED: description of top-ranked coherence pairs and
frequency bands; identification of most important spectral power
features; relative contribution of Riemannian covariance features.]}
Taken together, these findings indicate that the site effect in
ELM19 is primarily encoded in the inter-electrode coherence structure,
particularly in \textbf{[TO BE FILLED]} frequency bands, implicating
\textbf{[TO BE FILLED: procedural origin, e.g.\ differences in
eye-state protocols or patient preparation]}.
Figure~\ref{fig:shap_topo} shows SHAP topographies for the three most
discriminable institutions (\textbf{[TO BE FILLED: e.g., MOR, B2K,
TOR]}).
\textbf{[TO BE FILLED: describe dominant coherence patterns visible in
each.]}

\begin{figure}[htbp]
  \centering
  % \includegraphics[width=\textwidth]{Figs/shap_topographies.pdf}
  \caption{SHAP topographies for selected highly discriminable
    institutions, showing the spatial distribution of the most
    important features for site identification.}
  \label{fig:shap_topo}
\end{figure}

\begin{table}[htbp]
  \centering
  \caption{Overall site-classification MCC (mean $\pm$ SD, 5-fold CV
    on $\mathrm{ELM}_n$) for the feature-group ablation study.
    Upper section: models trained on a single feature group.
    Lower section: models trained with one group excluded.}
  \label{tab:ablation}
  \begin{tabular}{lcc}
    \toprule
    \textbf{Condition} & \textbf{MCC} & \textbf{SD} \\
    \midrule
    \multicolumn{3}{l}{\textit{Single feature group}} \\
    \quad Coherence only  & [TO BE FILLED] & [TO BE FILLED] \\
    \quad Power only      & [TO BE FILLED] & [TO BE FILLED] \\
    \quad Covariance only & [TO BE FILLED] & [TO BE FILLED] \\
    \midrule
    \multicolumn{3}{l}{\textit{One group excluded}} \\
    \quad All except coherence  & [TO BE FILLED] & [TO BE FILLED] \\
    \quad All except power      & [TO BE FILLED] & [TO BE FILLED] \\
    \quad All except covariance & [TO BE FILLED] & [TO BE FILLED] \\
    \bottomrule
  \end{tabular}
\end{table}

\subsection{Feature-Level Harmonisation: Univariate and Multivariate
Assessment}
\label{subsec:harmonisation_results}

\subsubsection{Cohen's~\texorpdfstring{$d$}{d}: Univariate
Harmonisation Effectiveness}

Prior to harmonisation, substantial inter-site differences were
present at the feature level, with mean pairwise Cohen's~$d$
of \textbf{[TO BE FILLED]} averaged across all institution pairs and
all 2\,850 features.
Table~\ref{tab:cohend_summary} reports this summary for all
harmonisation conditions.
All marginal correction methods reduced inter-site Cohen's~$d$
relative to the unharmonised baseline.
\textbf{[TO BE FILLED: whether PCA pre-orthogonalisation and/or CovBat
provide additional attenuation, and how the conditions rank.]}
This result confirms that all methods achieve their intended goal at
the univariate level.

Figure~\ref{fig:cohend_heatmap} shows pairwise inter-institution
Cohen's~$d$ before and after \textbf{[TO BE FILLED: method]}.
\textbf{[TO BE FILLED: describe pattern.]}

\begin{figure}[htbp]
  \centering
  % \includegraphics[width=\textwidth]{Figs/cohens_d_heatmap.pdf}
  \caption{Pairwise inter-institution mean $|$Cohen's~$d|$ across all
    2\,850 features, before (left) and after (right) harmonisation.
    Institutions are ordered identically in both panels.}
  \label{fig:cohend_heatmap}
\end{figure}

To directly assess the multivariate covariance structure, pairwise
Frobenius distances between per-institution covariance matrices were
computed before and after harmonisation
(Figure~\ref{fig:covariance_distances}).
\textbf{[TO BE FILLED: describe whether distances persist after
harmonisation, supporting the paradox mechanism.]}

\begin{figure}[htbp]
  \centering
  % \includegraphics[width=\textwidth]{Figs/frobenius_distance.png}
  \caption{Pairwise Frobenius distances between per-institution
    covariance matrices. Persistent inter-site covariance structure
    after marginal harmonisation supports the mechanistic basis of
    the harmonisation paradox.}
  \label{fig:covariance_distances}
\end{figure}

The crucial question---addressed in the following section---is whether
univariate correction corresponds to a reduction of multivariate site
discriminability.

\begin{table}[htbp]
  \centering
  \caption{Mean pairwise Cohen's~$d$ (averaged across all institution
    pairs and all 2\,850 features) before and after each harmonisation
    condition, applied to $\mathrm{ELM}_n$.}
  \label{tab:cohend_summary}
  \begin{tabular}{lc}
    \toprule
    \textbf{Condition} & \textbf{Mean pairwise Cohen's~$d$} \\
    \midrule
    Original (unharmonised)         & [TO BE FILLED] \\
    PCA pre-orthogonalisation only  & [TO BE FILLED] \\
    Site-wise standardisation       & [TO BE FILLED] \\
    PCA + site-wise stand.          & [TO BE FILLED] \\
    neuroCombat                     & [TO BE FILLED] \\
    PCA + neuroCombat               & [TO BE FILLED] \\
    CovBat                          & [TO BE FILLED] \\
    PCA + CovBat                    & [TO BE FILLED] \\
    \bottomrule
  \end{tabular}
\end{table}

\subsubsection{Feature-Space Visualisation}
\label{subsubsec:tsne}

Figure~\ref{fig:tsne} shows $t$-SNE projections of the
$\mathrm{ELM}_n$ feature space before and after harmonisation,
coloured by institution identity.
\textbf{[TO BE FILLED: description of clustering patterns---do
hospital clusters persist after harmonisation?]}

\begin{figure}[htbp]
  \centering
  \begin{subfigure}[t]{0.48\textwidth}
    \centering
    % \includegraphics[width=\textwidth]{Figs/raw_tsne.png}
    \caption{Raw (unharmonised)}
  \end{subfigure}
  \hfill
  \begin{subfigure}[t]{0.48\textwidth}
    \centering
    % \includegraphics[width=\textwidth]{Figs/sitewise_tsne.png}
    \caption{Site-wise standardisation}
  \end{subfigure}
  \\[6pt]
  \begin{subfigure}[t]{0.48\textwidth}
    \centering
    % \includegraphics[width=\textwidth]{Figs/neurocombat_tsne.png}
    \caption{neuroCombat}
  \end{subfigure}
  \hfill
  \begin{subfigure}[t]{0.48\textwidth}
    \centering
    % \includegraphics[width=\textwidth]{Figs/covbat_tsne.png}
    \caption{CovBat}
  \end{subfigure}
  \caption{$t$-SNE projections of the $\mathrm{ELM}_n$ feature space
    coloured by institution identity, before and after harmonisation.
    \textbf{[TO BE FILLED: describe persistence of clusters.]}}
  \label{fig:tsne}
\end{figure}

\subsubsection{The Harmonisation Paradox: Multivariate Site
Discriminability After Correction}
\label{subsubsec:paradox}

Figure~\ref{fig:site_mcc_no_pca} shows the per-institution
site-classification MCC for all harmonisation methods applied to the
original, non-orthogonalised feature space.
Without PCA pre-orthogonalisation, all correction methods drive
per-site MCC towards a near-ceiling value close to 1.0 across
virtually all institutions---a uniformly strong manifestation of the
\emph{harmonisation paradox}.
The overall site-classification MCC for each condition is summarised
in Table~\ref{tab:site_class_harmonisation}.
This result demonstrates that marginal correction, while reducing
univariate inter-site differences (Table~\ref{tab:cohend_summary}),
simultaneously reorganises the multivariate feature space in a way
that makes the recording institution \emph{more} identifiable, not
less.
In contrast, logistic regression applied to the same harmonised
features showed \textbf{[TO BE FILLED: reduced/unchanged]}
site-classification MCC after marginal correction, indicating that
the paradox is specific to classifiers with sufficient capacity to
exploit inter-feature covariance structure.
The mechanism is the following: by equalising feature means and
variances, marginal correction removes the dominant sources of
within-site variance, causing the subtler but highly structured
inter-site covariance patterns---which survive the correction
intact---to become more salient to the nonlinear classifier.

\begin{figure}[htbp]
  \centering
  \includegraphics[width=\textwidth]{Figs/site_mcc_pca_00.pdf}
  \caption{Per-institution site-classification MCC (one-vs-rest,
    5-fold CV on $\mathrm{ELM}_n$) for the raw feature set and after
    applying four harmonisation methods \emph{without} PCA
    pre-orthogonalisation.
    Institutions are ordered by decreasing raw-baseline MCC.
    All four methods drive MCC towards a near-ceiling value,
    demonstrating the harmonisation paradox.}
  \label{fig:site_mcc_no_pca}
\end{figure}

Figure~\ref{fig:site_mcc_pca} shows the analogous results when PCA
pre-orthogonalisation (retaining 100\% of variance) is applied to
the feature space prior to harmonisation.
PCA pre-orthogonalisation substantially attenuates the paradox: the
raw-baseline MCC is markedly reduced across institutions, and the
harmonised conditions no longer approach ceiling performance.
Instead, per-site MCC values show high institution-level variability,
with harmonised conditions fluctuating around---and in several cases
falling below---the raw baseline.
This demonstrates that decorrelating the feature space prior to
marginal correction eliminates the covariance structure that drives
the paradox, consistent with the mechanistic explanation above.

\begin{figure}[htbp]
  \centering
  \includegraphics[width=\textwidth]{Figs/site_mcc_pca_10.pdf}
  \caption{Per-institution site-classification MCC (one-vs-rest,
    5-fold CV on $\mathrm{ELM}_n$) for the raw feature set and after
    applying four harmonisation methods \emph{with} full-variance PCA
    pre-orthogonalisation.
    Institutions are ordered by decreasing raw-baseline MCC.
    PCA pre-orthogonalisation substantially attenuates the
    harmonisation paradox.}
  \label{fig:site_mcc_pca}
\end{figure}

\begin{table}[htbp]
  \centering
  \caption{Overall site-classification MCC (mean $\pm$ SD, 5-fold CV
    on $\mathrm{ELM}_n$) for all harmonisation conditions.
    Values above the unharmonised baseline quantify the degree to
    which the harmonisation paradox persists under each condition.
    The bottom section reports linear-probe MCC on the recording-level
    representations $\mathbf{h}$ from the neural models.}
  \label{tab:site_class_harmonisation}
  \begin{tabular}{lcc}
    \toprule
    \textbf{Condition} & \textbf{MCC} & \textbf{SD} \\
    \midrule
    \multicolumn{3}{l}{\textit{CatBoost, without PCA}} \\
    \quad Original (unharmonised)        & [TO BE FILLED] & [TO BE FILLED] \\
    \quad Site-wise standardisation      & [TO BE FILLED] & [TO BE FILLED] \\
    \quad neuroCombat                    & [TO BE FILLED] & [TO BE FILLED] \\
    \quad CovBat                         & [TO BE FILLED] & [TO BE FILLED] \\
    \midrule
    \multicolumn{3}{l}{\textit{Logistic regression, without PCA}} \\
    \quad Original (unharmonised)        & [TO BE FILLED] & [TO BE FILLED] \\
    \quad Site-wise standardisation      & [TO BE FILLED] & [TO BE FILLED] \\
    \quad neuroCombat                    & [TO BE FILLED] & [TO BE FILLED] \\
    \quad CovBat                         & [TO BE FILLED] & [TO BE FILLED] \\
    \midrule
    \multicolumn{3}{l}{\textit{CatBoost, with PCA pre-orthogonalisation}} \\
    \quad PCA only                       & [TO BE FILLED] & [TO BE FILLED] \\
    \quad PCA + site-wise stand.         & [TO BE FILLED] & [TO BE FILLED] \\
    \quad PCA + neuroCombat              & [TO BE FILLED] & [TO BE FILLED] \\
    \quad PCA + CovBat                   & [TO BE FILLED] & [TO BE FILLED] \\
    \midrule
    % TODO(Julia): Fill MINet/MINet-DANN results
    \multicolumn{3}{l}{\textit{Signal-level (linear probe on $\mathbf{h}$)}} \\
    \quad MINet (no adversarial training) & [TO BE FILLED] & [TO BE FILLED] \\
    \quad MINet-DANN                      & [TO BE FILLED] & [TO BE FILLED] \\
    \bottomrule
  \end{tabular}
\end{table}

\subsection{PCA Variance Sensitivity Analysis}
\label{subsec:pca_sweep}

\textbf{[TO BE FILLED: table or figure showing site MCC $\times$
method $\times$ PCA variance level.]}

\textbf{[TO BE FILLED: table or figure showing pathology AUC
$\times$ method $\times$ PCA variance level.]}

Retaining 90\% of variance yielded the best trade-off:
site-classification MCC was reduced by \textbf{[TO BE FILLED]}
relative to full variance, while pathology AUC was
\textbf{[TO BE FILLED: maintained/improved]} relative to the
unharmonised baseline.

Figure~\ref{fig:tradeoff} visualises the per-institution trade-off
between site-invariance ($\Delta$Site MCC, $x$-axis) and clinical
utility ($\Delta$Pathology AUC, $y$-axis) for each harmonisation
method across PCA variance levels.
\textbf{[TO BE FILLED: describe the quadrant shift from bottom-right
to top-left with PCA.]}

\begin{figure}[htbp]
  \centering
  % \includegraphics[width=\textwidth]{Figs/tradeoff_grid.png}
  \caption{Per-institution trade-off between site-invariance
    ($\Delta$Site MCC) and clinical utility ($\Delta$Pathology AUC)
    for each harmonisation method across PCA variance levels.}
  \label{fig:tradeoff}
\end{figure}

\subsection{Robustness of the Site Effect to Preprocessing}
\label{subsec:filter_robustness}

The site-classification MCC remained stable across all five
preprocessing configurations, ranging from \textbf{[TO BE FILLED]}
to \textbf{[TO BE FILLED]}
(Table~\ref{tab:filter_robustness}), confirming that the site effect
is not an artefact of a specific filter choice.

\begin{table}[htbp]
  \centering
  \caption{Site-classification MCC (mean $\pm$ SD, 5-fold CV on
    $\mathrm{ELM}_n$) across five preprocessing filter
    configurations.}
  \label{tab:filter_robustness}
  \begin{tabular}{lcc}
    \toprule
    \textbf{Filter Configuration} & \textbf{MCC} & \textbf{SD} \\
    \midrule
    \textbf{[TO BE FILLED from
      results/tables/06\_filter\_robustness/site\_clf\_results.csv]} \\
    \bottomrule
  \end{tabular}
\end{table}

\begin{figure}[htbp]
  \centering
  % \includegraphics[width=\textwidth]{Figs/site_mcc_overall_bar.png}
  \caption{Overall site-classification MCC across five preprocessing
    filter configurations.}
  \label{fig:filter_robustness_bar}
\end{figure}

\subsection{Impact of Harmonisation on Clinical Neuroscreening}
\label{subsec:clinical_results}

\subsubsection{Per-Institution LOSO Performance}

Figure~\ref{fig:patho_auc_no_pca} shows the per-institution
neuroscreening AUC under the LOSO protocol for all harmonisation
methods applied without PCA pre-orthogonalisation.
Without PCA, the four methods diverge noticeably at individual
institutions, with CovBat and neuroCombat showing anomalous behaviour
at specific sites.
\textbf{[TO BE FILLED: identify the specific outlier institutions,
describe the direction of the anomaly, and comment on likely causes
(e.g.\ small calibration sample, specific demographic or protocol
characteristics).]}
The ordering of institutions by raw-baseline AUC broadly follows the
pattern observed in the site-classification MCC, with the
lowest-performing sites---notably PIO and TER\_L---showing near-chance
performance for all methods, suggesting that these institutions
present an inherently difficult generalisation challenge irrespective
of harmonisation.

\begin{figure}[htbp]
  \centering
  \includegraphics[width=\textwidth]{Figs/patho_auc_pca_none.pdf}
  \caption{Per-institution neuroscreening AUC under the LOSO protocol
    for the GBE classifier applied to raw features and after four
    harmonisation methods, \emph{without} PCA pre-orthogonalisation.
    Institutions are ordered by decreasing raw-baseline AUC.}
  \label{fig:patho_auc_no_pca}
\end{figure}

Figure~\ref{fig:patho_auc_pca} shows the analogous results when PCA
pre-orthogonalisation is applied prior to harmonisation.
In contrast to the no-PCA condition, the four methods cluster closely
together across all institutions, and the institution-level
variability in AUC is markedly reduced.
The harmonised conditions consistently track slightly above the raw
baseline for the majority of institutions, indicating that the
combination of PCA pre-orthogonalisation with any of the tested
marginal correction methods yields a more uniform and stable
improvement in cross-site neuroscreening performance than marginal
correction alone.
\textbf{[TO BE FILLED: identify any remaining outlier institutions
and comment on whether the stabilising effect of PCA is uniform
across the full performance range or more pronounced at specific sites.]}

\begin{figure}[htbp]
  \centering
  \includegraphics[width=\textwidth]{Figs/patho_auc_pca_all.pdf}
  \caption{Per-institution neuroscreening AUC under the LOSO protocol
    for the GBE classifier applied to raw features and after four
    harmonisation methods, \emph{with} full-variance PCA
    pre-orthogonalisation.
    Institutions are ordered by decreasing raw-baseline AUC.
    PCA pre-orthogonalisation stabilises the harmonisation effect:
    all four methods cluster closely and track slightly above the raw
    baseline, with substantially reduced institution-level variability
    compared to the no-PCA condition
    (Figure~\ref{fig:patho_auc_no_pca}).}
  \label{fig:patho_auc_pca}
\end{figure}

\subsubsection{Aggregate Neuroscreening Performance}

Tables~\ref{tab:clinical_cv} and~\ref{tab:clinical_loso} report
aggregate neuroscreening MCC and AUC under 5-fold cross-validation
and LOSO, respectively, for all conditions including the MINet
baseline and MINet-DANN.
\textbf{[TO BE FILLED: narrative summary identifying whether
harmonisation consistently improves on the unharmonised baseline,
whether PCA pre-orthogonalisation provides additional aggregate
benefit, how CovBat compares to neuroCombat and site-wise
standardisation, and where MINet-DANN sits relative to feature-level
methods.]}

\begin{table}[htbp]
  \centering
  \caption{Neuroscreening performance under 5-fold cross-validation
    (mean $\pm$ SD) for all conditions.}
  \label{tab:clinical_cv}
  \begin{tabular}{lcc}
    \toprule
    \textbf{Method} & \textbf{MCC} & \textbf{AUC} \\
    \midrule
    \multicolumn{3}{l}{\textit{GBE, without PCA pre-orthogonalisation}} \\
    \quad Unharmonised               & [TO BE FILLED] & [TO BE FILLED] \\
    \quad Site-wise standardisation  & [TO BE FILLED] & [TO BE FILLED] \\
    \quad neuroCombat                & [TO BE FILLED] & [TO BE FILLED] \\
    \quad CovBat                     & [TO BE FILLED] & [TO BE FILLED] \\
    \midrule
    \multicolumn{3}{l}{\textit{GBE, with PCA pre-orthogonalisation}} \\
    \quad PCA only                   & [TO BE FILLED] & [TO BE FILLED] \\
    \quad PCA + site-wise stand.     & [TO BE FILLED] & [TO BE FILLED] \\
    \quad PCA + neuroCombat          & [TO BE FILLED] & [TO BE FILLED] \\
    \quad PCA + CovBat               & [TO BE FILLED] & [TO BE FILLED] \\
    \midrule
    % TODO(Julia): Fill MINet/MINet-DANN results
    \multicolumn{3}{l}{\textit{Signal-level adversarial harmonisation}} \\
    \quad MINet (no adv.\ training)  & [TO BE FILLED] & [TO BE FILLED] \\
    \quad MINet-DANN                 & [TO BE FILLED] & [TO BE FILLED] \\
    \bottomrule
  \end{tabular}
\end{table}

\begin{table}[htbp]
  \centering
  \caption{Neuroscreening performance under the Leave-One-Site-Out
    protocol (mean $\pm$ SD) for all conditions.}
  \label{tab:clinical_loso}
  \begin{tabular}{lcc}
    \toprule
    \textbf{Method} & \textbf{MCC} & \textbf{AUC} \\
    \midrule
    \multicolumn{3}{l}{\textit{GBE, without PCA pre-orthogonalisation}} \\
    \quad Unharmonised               & [TO BE FILLED] & [TO BE FILLED] \\
    \quad Site-wise standardisation  & [TO BE FILLED] & [TO BE FILLED] \\
    \quad neuroCombat                & [TO BE FILLED] & [TO BE FILLED] \\
    \quad CovBat                     & [TO BE FILLED] & [TO BE FILLED] \\
    \midrule
    \multicolumn{3}{l}{\textit{GBE, with PCA pre-orthogonalisation}} \\
    \quad PCA only                   & [TO BE FILLED] & [TO BE FILLED] \\
    \quad PCA + site-wise stand.     & [TO BE FILLED] & [TO BE FILLED] \\
    \quad PCA + neuroCombat          & [TO BE FILLED] & [TO BE FILLED] \\
    \quad PCA + CovBat               & [TO BE FILLED] & [TO BE FILLED] \\
    \midrule
    % TODO(Julia): Fill MINet/MINet-DANN results
    \multicolumn{3}{l}{\textit{Signal-level adversarial harmonisation}} \\
    \quad MINet (no adv.\ training)  & [TO BE FILLED] & [TO BE FILLED] \\
    \quad MINet-DANN                 & [TO BE FILLED] & [TO BE FILLED] \\
    \bottomrule
  \end{tabular}
\end{table}

% -----------------------------------------------------------------------
\section{Discussion}
\label{sec:discussion}

\subsection{The Nature and Magnitude of the Site Effect}
\label{subsec:disc_site_effect}

The site-classification experiment confirmed that the site effect is
both strong and broadly distributed across the 2\,850-dimensional
feature space of the ELM19 corpus, with an overall MCC of
$\mathbf{[TO\ BE\ FILLED]}$.
\textbf{[TO BE FILLED: discussion of the range of per-site MCC,
identification of the most and least discriminable institutions, and
interpretation in terms of their demographic and protocol
characteristics as established in Section~\ref{subsec:exploratory}.]}

The SHAP analysis and ablation study converge on a consistent
interpretation: inter-electrode coherence features are the primary
carriers of site-specific information.
The elevated importance of O1--O2 coherence is consistent with
the hypothesis that differences in eye-open versus eye-closed
protocols across institutions modulate occipital alpha coherence,
which is highly sensitive to visual state.
The prominence of beta-band coherence in temporal regions is
consistent with variability in muscle artefact levels arising from
differences in patient preparation, scalp cleaning, and electrode
impedance management across sites, since the beta band overlaps
with the electromyographic spectrum.
Delta-band power and frontal features, identified as important
spectral components, may reflect differences in patient alertness
and vigilance state at recording time---sources of inter-site
variability that are difficult to control retrospectively.

The ELM19 corpus (54\,779 recordings, 30 institutions) is
substantially larger than previously studied multi-site EEG datasets,
including HarMNqEEG (1\,564 subjects, 9
countries)~\cite{li2022harmnqeeg} and the five-dataset pool of
Jaramillo-Jimenez et~al.~\cite{jaramillo2024combat}.

The robustness analysis across five preprocessing filter
configurations further confirms that the site effect reflects genuine
inter-site differences rather than an artefact of a specific filter
choice (Table~\ref{tab:filter_robustness}).

\subsection{The Harmonisation Paradox and Its Mechanistic Basis}
\label{subsec:disc_paradox}

The most striking finding of this study is the \emph{harmonisation
paradox}: all marginal correction methods successfully reduced
pairwise Cohen's~$d$ inter-site differences at the feature level,
yet simultaneously increased the site-classification MCC towards
a near-ceiling value.
This result, which at first appears contradictory, is consistent with
the theoretical limitations of feature-by-feature correction.

Marginal correction adjusts the mean and variance of each feature
independently, thereby equalising marginal distributions across
sites.
However, the site effect is encoded not only in the marginal
distribution of individual features but also in the
\emph{covariance structure} between features---a form of inter-site
variability that feature-by-feature corrections leave entirely
intact~\cite{hu2023image}.
By equalising means and variances, these methods remove the dominant
sources of within-site variance, causing the subtler but
highly structured inter-feature covariance patterns that encode site
identity to become more salient to the classifier.
This argument is analogous to that advanced by
Mellot~et~al.~\cite{mellot2023meeg} in the context of Riemannian
geometry-based EEG alignment: corrections that ignore covariance
structure are fundamentally insufficient to eliminate multivariate
site effects.

The PCA pre-orthogonalisation experiments provide direct mechanistic
evidence for this interpretation.
By projecting the feature space onto mutually uncorrelated principal
components, PCA destroys the inter-feature covariance structure that
drives the paradox.
Marginal correction applied to the orthogonalised space therefore
cannot amplify covariance-based site signatures, and the paradox is
substantially attenuated.
The classifier-dependence of the paradox provides additional
mechanistic evidence.
Logistic regression, which models linear decision boundaries and
cannot exploit inter-feature covariance, does not exhibit the
paradox: harmonisation successfully reduces site discriminability
under a linear classifier.
CatBoost, by contrast, learns nonlinear decision surfaces that can
detect site-specific patterns in the residual covariance structure.
This demonstrates that the paradox is not a property of the data
alone but of the interaction between residual multivariate structure
and classifier expressiveness.

The PCA variance sweep reveals that retaining 90\% of variance
provides the optimal trade-off between site-invariance and clinical
utility.
\textbf{[TO BE FILLED: interpret why removing 10\% of variance
helps.]}

Importantly, PCA pre-orthogonalisation also stabilises the
institution-level variability of neuroscreening performance across
sites (Figures~\ref{fig:patho_auc_no_pca}
and~\ref{fig:patho_auc_pca}), suggesting that covariance-based
site structure contributes not only to the paradox but also to the
institution-specific performance fluctuations observed without PCA.

This finding has a direct practical implication for harmonisation
methodology: univariate metrics such as Cohen's~$d$ are inadequate
as sole indicators of harmonisation quality in high-dimensional
feature spaces.
A complete evaluation framework must include a multivariate
site-discriminability measure---such as the site-classification MCC
reported here---as a necessary complement to feature-level
effect-size summaries.

\subsection{Impact on Clinical Neuroscreening}
\label{subsec:disc_clinical}

Harmonisation consistently \textbf{[TO BE FILLED: improved/maintained]}
pathology classification under both CV and LOSO protocols.
\textbf{[TO BE FILLED: which method performed best.]}
PCA pre-orthogonalisation \textbf{[TO BE FILLED: effect on pathology
results.]}
The LOSO results are particularly important because they directly
simulate deployment to a previously unseen site, the most clinically
relevant scenario.
\textbf{[TO BE FILLED: whether the LOSO and CV results are concordant
or reveal differences; where MINet-DANN sits relative to feature-level
conditions.]}

\subsection{Limitations}
\label{subsec:limitations}

Several limitations of the present study should be noted.
First, although patient age and sex were included as biological
covariates in neuroCombat and CovBat, other clinically relevant
variables---such as the reason for referral, medication status, or
cognitive state at recording time---were not available for the full
ELM19 corpus and could not be controlled for; residual confounding
from these sources may inflate apparent site differences.
Second, all marginal correction methods adjust only the first two
moments of each feature's distribution and cannot correct
higher-order distributional differences (e.g.\ skewness,
multimodality) that may contribute to the site effect.
Third, the LOSO calibration procedure (using $N = 30$ normal
recordings from the held-out site) introduces a practical dependency
on the availability of a small representative sample from any new
site, which may not be feasible in all prospective deployments.
Fourth, PCA pre-orthogonalisation discards the covariance structure
of the original feature space; while this attenuates the harmonisation
paradox, it also removes inter-feature relationships that may carry
clinically relevant information, and the long-term impact on
downstream task performance warrants further investigation.
Fifth, the MINet-DANN model was evaluated on the same corpus as the
feature-level methods, but a direct comparison under identical
conditions across all 30 LOSO rounds requires substantial
computational resources and results may not yet be complete at the
time of submission \textbf{[DANN --- Julia to update]}.
Sixth, the logistic regression comparison was conducted with only two
classifier families; the classifier-dependence of the paradox may
manifest differently with other model families (e.g.\ random forests,
SVMs) and warrants further investigation.

% TODO(Julia): Update this section based on DANN results
\subsection{Future Directions}
\label{subsec:future}

The harmonisation paradox and its mechanistic basis identified in
this study motivate several directions for future work.
Most immediately, the comparison between MINet-DANN and the
feature-level methods provides a first assessment of whether
end-to-end adversarial harmonisation can overcome the fundamental
limitation of marginal correction.
The linear probe site-classification results on the latent
representations $\mathbf{h}$ of MINet vs.\ MINet-DANN
(Table~\ref{tab:site_class_harmonisation}) provide a direct parallel
to the feature-level site-classification analysis and will determine
whether adversarial training successfully suppresses multivariate
site information at the representation level.

Beyond harmonisation methodology, the present results also point
toward the value of feature engineering strategies that reduce
site sensitivity by design.
The ablation study and SHAP analysis suggest that power features carry
substantially less site-specific information than coherence features,
indicating that feature selection guided by SHAP importance could
reduce the effective site effect without explicit harmonisation.
Finally, normative modelling frameworks~\cite{dong2025normative},
which express individual measurements as deviations from an
age-adjusted population norm, constitute a complementary approach that
may interact favourably with the harmonisation methods studied here
and warrant systematic evaluation in the ELM19 context.

% -----------------------------------------------------------------------
\section{Conclusion}
\label{sec:conclusion}

This study provides a comprehensive characterisation of the site
effect in the ELM19 large-scale clinical EEG corpus and a systematic
evaluation of feature-level harmonisation strategies including
site-wise standardisation, neuroCombat, CovBat, and PCA
pre-orthogonalisation, complemented by a signal-level adversarial
baseline.
A gradient-boosted classifier achieved a site-identification MCC of
$\mathbf{[TO\ BE\ FILLED]}$, confirming that site identity is
strongly encoded in EEG feature representations.
SHAP explainability analysis and ablation experiments identified
inter-electrode coherence---particularly O1--O2 and beta-band
temporal coherence---as the dominant carrier of site-specific
information, implicating differences in eye-state protocols and
muscle artefact management across institutions.

A central finding is the \emph{harmonisation paradox}: all marginal
correction methods reduced univariate inter-site differences
(Cohen's~$d$) yet simultaneously increased the site-classification
MCC to a near-ceiling value, demonstrating that feature-by-feature
correction reorganises rather than eliminates the multivariate
covariance structure of the site effect.
PCA pre-orthogonalisation of the feature space before correction
substantially attenuates this paradox and yields more uniform and
stable cross-site neuroscreening performance across institutions.
A sensitivity analysis across PCA retained-variance levels established
that retaining 90\% of variance provides the optimal trade-off
between site-invariance and clinical utility.
The paradox was shown to be classifier-dependent: logistic regression
did not exhibit the paradox, confirming that it arises from the
interaction between residual covariance structure and nonlinear
classifier capacity.
The site effect was further shown to be robust across five
preprocessing filter configurations, confirming that it reflects
genuine inter-site differences rather than preprocessing artefacts.
These findings demonstrate that univariate metrics alone are
insufficient to validate harmonisation quality in high-dimensional
feature spaces, and that evaluation frameworks must include a
multivariate site-discriminability measure.

% TODO(Julia): Add DANN conclusion summary
\textbf{[TO BE FILLED: summary of aggregate neuroscreening results
and their interpretation once all conditions are complete.]}

These results collectively underscore that the site effect is a
complex, high-dimensional challenge for which standard marginal
correction methods provide only a partial solution.
Developing truly generalisable neuroscreening models will require
harmonisation strategies that explicitly account for multivariate
inter-site structure---whether through covariance-aware statistical
corrections, adversarial representation learning, or feature
engineering guided by explainability analysis.

% -----------------------------------------------------------------------
\section*{Acknowledgements}
The authors are grateful to Elmiko Biosignals sp.~z~o.o.\ for
providing access to the ELM19 database.
% [Add funding acknowledgements here.]

\section*{CRediT Author Contribution Statement}
K.C.: Methodology, Software, Formal Analysis, Investigation,
Visualisation, Writing -- Original Draft.
J.P.: Methodology, Writing -- Review \& Editing.
J.\.Z.: Conceptualisation, Supervision, Resources, Writing --
Review \& Editing.

\section*{Declaration of Competing Interests}
The authors declare that they have no known competing financial
interests or personal relationships that could have appeared to
influence the work reported in this paper.

\section*{Data Availability}
The ELM19 database is available from Elmiko Biosignals sp.~z~o.o.\
upon reasonable request to
\href{mailto:EEG.database@elmiko.pl}{EEG.database@elmiko.pl}.

% -----------------------------------------------------------------------
\appendix
\section{Patient Demographics by Institution}
\label{app:demographics}

\begin{table}[htbp]
  \centering
  \caption{Summary statistics of patient age (years) by institution
    for the $\mathrm{ELM}_n$ normal recording subset.
    A representative selection of 9 institutions is shown here;
    the complete table for all 30 institutions is available in the
    supplementary materials.}
  \label{tab:age_stats}
  \begin{tabular}{lrrrrr}
    \toprule
    \textbf{Institution} & \textbf{Mean} & \textbf{Median} &
    \textbf{SD} & \textbf{Min} & \textbf{Max} \\
    \midrule
    PIO    & 20.3 & 17.6 & 8.3  & 16.0 & 68.3 \\
    SL2    & 28.4 & 19.5 & 15.1 & 16.0 & 66.3 \\
    TER\_L & 26.3 & 17.8 & 15.0 & 16.0 & 69.8 \\
    MKW    & 26.1 & 21.7 & 10.8 & 16.0 & 69.7 \\
    GAK    & 41.3 & 40.9 & 13.8 & 16.0 & 70.0 \\
    KAL    & 46.2 & 48.2 & 15.6 & 17.3 & 69.9 \\
    OST    & 48.1 & 49.2 & 13.7 & 18.4 & 69.9 \\
    MOR    & 48.6 & 52.1 & 15.4 & 18.7 & 69.4 \\
    TOR    & 47.6 & 50.8 & 16.5 & 16.0 & 69.9 \\
    \multicolumn{6}{c}{$\cdots$} \\
    \bottomrule
  \end{tabular}
\end{table}

% -----------------------------------------------------------------------
\bibliographystyle{elsarticle-num}
\bibliography{bibliography}

\end{document}